\documentclass[11pt]{article}

% ---------- Packages ----------
\usepackage[margin=1in]{geometry}
\usepackage{microtype}
\usepackage{booktabs}
\usepackage{multirow}
\usepackage{graphicx}
\usepackage{subcaption}
\usepackage{amsmath, amssymb, amsthm}
\usepackage{mathtools}
\usepackage{algorithm}
\usepackage{algpseudocode}
\usepackage{hyperref}

% ---------- Theorems ----------
\newtheorem{theorem}{Theorem}
\newtheorem{lemma}{Lemma}
\newtheorem{proposition}{Proposition}
\newtheorem{corollary}{Corollary}
\theoremstyle{definition}
\newtheorem{definition}{Definition}

% ---------- Notation ----------
\newcommand{\R}{\mathbb{R}}
\newcommand{\Cl}{\mathrm{Cl}}
\newcommand{\SO}{\mathrm{SO}}
\newcommand{\Spin}{\mathrm{Spin}}
\newcommand{\so}{\mathfrak{so}}
\newcommand{\Ad}{\mathrm{Ad}}
\newcommand{\tr}{\mathrm{tr}}
\newcommand{\norm}[1]{\left\lVert #1 \right\rVert}
\newcommand{\inner}[2]{\left\langle #1, #2 \right\rangle}

% Skew-sym / commutator
\newcommand{\skew}[1]{\mathrm{skew}\!\left(#1\right)}
\newcommand{\comm}[2]{\left[ #1, #2 \right]}

% Cayley
\newcommand{\Cay}{\mathrm{Cay}}


\title{\textbf{BNR-PE: Budgeted Non-Abelian Rotor Positional Encodings via Geometric Algebra}}
\author{Anonymous Authors}
\date{}

\begin{document}
\maketitle

\begin{abstract}
Rotary positional encodings (RoPE) implement positional structure via commuting 2D rotations, yielding exact relative-offset equivariance and excellent long-context extrapolation but limited cross-dimensional coupling. Lie relative encodings (LieRE) generalize RoPE by learning dense skew-symmetric generators and exponentiating them, increasing expressivity for multi-dimensional domains at the cost of heavier computation and loss of exact relative structure. We introduce \emph{Budgeted Non-Abelian Rotor Positional Encodings (BNR-PE)}: a geometric-algebra formulation that treats skew generators as bivectors and constructs rotors with \emph{structured non-commutativity} controlled by an explicit \emph{commutator budget}. BNR-PE interpolates between RoPE (Abelian) and LieRE (non-Abelian) while preserving norm (hypersphere) compatibility and enabling efficient implementations via (i) low-rank and block-sparse bivector parameterizations, (ii) a Cayley--Woodbury rotor map that avoids dense matrix exponentials while remaining exactly orthogonal, and (iii) differential rotor accumulation for long sequences and grids. We provide bounds quantifying deviation from pure relative-offset behavior under bounded commutators and validate BNR-PE on text, 2D/3D, and time--frequency tasks with single-digit overhead over FlashAttention-style baselines.
\end{abstract}

\section{Introduction}
We seek positional encodings that (i) preserve RoPE's strong long-context extrapolation, (ii) generalize cleanly to multi-dimensional coordinates (images, video, time--frequency), and (iii) admit kernel-efficient implementations compatible with modern fused attention backends.
RoPE achieves (i) by restricting to commuting 2D rotations, but sacrifices cross-dimensional coupling.
LieRE increases expressivity via learned dense rotations but introduces additional compute and degrades exact relative-offset structure.

\paragraph{Thesis.}
We propose \emph{Budgeted Non-Abelian Rotor Positional Encodings (BNR-PE)}: a rotor-based scheme whose non-commutativity is \emph{structured} and \emph{budgeted} to balance expressivity and extrapolation.
Our primary technical device is a low-rank / block-sparse generator family and an exactly-orthogonal \emph{Cayley--Woodbury} rotor map that avoids dense matrix exponentials while retaining norm preservation.

\paragraph{Contributions.}
\begin{enumerate}
  \item A geometric-algebra (Clifford) view that treats skew generators as bivectors and unifies RoPE and LieRE in a single rotor framework.
  \item A commutator-budgeted generator class that interpolates between Abelian (RoPE) and non-Abelian (LieRE-like) behavior.
  \item Efficient rotor application via Cayley transforms plus Woodbury identities for low-rank generators, enabling near-RoPE overhead.
  \item A theory section establishing exact $\Delta$-only structure in the Abelian case and quantitative deviation bounds under bounded commutators.
  \item Empirical results and ablations across text, 2D/3D, and time--frequency domains.
\end{enumerate}


\section{Background}
\subsection{RoPE as commuting rotors}
RoPE partitions the embedding into $d/2$ orthogonal $2$-planes and applies independent rotations with position-dependent phases.
This yields the key identity $\widetilde{R(i)} R(j) = R(j-i)$, so attention scores depend only on relative offsets.

\subsection{LieRE as learned Lie algebra actions}
LieRE generalizes rotary encodings by learning skew-symmetric generators from multi-dimensional coordinates and exponentiating them to rotations.
The resulting rotations are typically non-commuting, producing richer structure at the cost of exponentiation overhead and loss of exact relative-only behavior.

\subsection{Geometric algebra preliminaries}
Let $V\cong\R^d$. The Clifford algebra $\Cl(V)$ provides a product extending the inner product on $V$.
Bivectors correspond to oriented $2$-planes and generate rotations (rotors) via exponentiation.
In this paper we use geometric algebra primarily as a unifying and constructive language; implementation reduces to structured skew-symmetric operators and orthogonal transforms.


\section{BNR-PE}
\subsection{Rotor positional encodings}
Let positions be $P\in\R^n$.
We define a skew-symmetric generator $G(P)\in\so(d)$ and the rotor (orthogonal map)
\begin{equation}
  R(P) \;\coloneqq\; \exp\!\left(-\tfrac12 G(P)\right)\in \SO(d).
\end{equation}
Given an input vector $x\in\R^d$ (query, key, or value), we apply the positional action $x \mapsto R(P)\,x$.

\subsection{BNR-PE generators with commutator budgets}
We use a linear coordinate map
\begin{equation}
  G(P) \;=\; \sum_{a=1}^{n} P_a\,A_a, \qquad A_a \in \so(d).
\end{equation}
Non-commutativity arises when $\comm{A_a}{A_b}\neq 0$.
We introduce an explicit \emph{commutator budget} $\kappa$ by constraining or regularizing
\begin{equation}
  \norm{\comm{A_a}{A_b}} \;\le\; \kappa \quad \forall a<b,
\end{equation}
so that the model interpolates between Abelian (RoPE-like) and non-Abelian (LieRE-like) regimes.

\subsection{Structured low-rank parameterization}
To enable efficient application, we parameterize each axis generator as a low-rank skew operator:
\begin{equation}
  A_a \;=\; U_a\,M_a\,U_a^\top,\qquad U_a\in\R^{d\times r}, \quad M_a^\top=-M_a \in \R^{r\times r},\quad r\ll d.
\end{equation}
This yields rank at most $r$ per axis and supports fast application via Woodbury identities.

\subsection{Cayley--Woodbury rotor map}
Instead of a dense exponential, we use the Cayley transform
\begin{equation}
  \Cay(S) \;\coloneqq\; (I-\tfrac12 S)^{-1}(I+\tfrac12 S),
  \qquad S^\top=-S,
\end{equation}
which is exactly orthogonal for skew $S$ and admits low-rank fast application when $S=UCU^\top$.
We define
\begin{equation}
  R(P) \;\coloneqq\; \Cay\!\big(G(P)\big),
\end{equation}
and study approximation error relative to $\exp(-\tfrac12 G(P))$ in Section~\ref{sec:theory}.


\section{Theory}
\label{sec:theory}
\subsection{Exact relative-offset structure in the Abelian case}
\begin{theorem}[Exact $\Delta$-only structure]
If $\comm{A_a}{A_b}=0$ for all $a,b$, then for all positions $P,Q\in\R^n$,
\begin{equation}
  R(P)^{-1}R(Q) \;=\; R(Q-P),
\end{equation}
and consequently attention logits constructed from dot products of rotated queries and keys depend only on $(Q-P)$.
\end{theorem}

\begin{proof}[Proof sketch]
When generators commute, $\exp(\cdot)$ and $\Cay(\cdot)$ both respect additive composition on the commuting subalgebra.
Thus $G(Q)-G(P)=G(Q-P)$ and the group product reduces to relative differences.
A full proof unrolls the commuting functional calculus and the identity $\Cay(S)^{-1}=\Cay(-S)$.
\end{proof}

\subsection{Deviation bounds under bounded commutators}
Let $X=-\tfrac12 G(P)$ and $Y=-\tfrac12 G(Q)$.
The Baker--Campbell--Hausdorff (BCH) formula gives
\begin{equation}
  \log\!\big(\exp(-X)\exp(Y)\big)
  \;=\;
  (Y-X) + \tfrac12\comm{Y}{X} + \tfrac1{12}\comm{Y}{\comm{Y}{X}} + \tfrac1{12}\comm{X}{\comm{X}{Y}} + \cdots.
\end{equation}
We bound $\norm{\comm{X}{Y}}$ in terms of the commutator budget.
\begin{lemma}[Commutator magnitude bound]
Assume $\norm{\comm{A_a}{A_b}}\le \kappa$ for all $a<b$.
Then
\begin{equation}
  \norm{\comm{X}{Y}}
  \;\le\;
  \frac{\kappa}{4}\sum_{a<b} \left| P_aQ_b - P_bQ_a \right|.
\end{equation}
\end{lemma}

\begin{proof}
By bilinearity,
$\comm{G(P)}{G(Q)}=\sum_{a,b}P_aQ_b\comm{A_a}{A_b}$.
Skew symmetry implies cancellation of diagonal terms; grouping $a<b$ yields coefficients $(P_aQ_b-P_bQ_a)$.
Apply the triangle inequality and the budget bound.
\end{proof}

\paragraph{Remark.}
The wedge-like term $P_aQ_b-P_bQ_a$ measures oriented ``area'' spanned by $(P,Q)$ in the $(a,b)$ coordinate plane.
Thus deviation from pure relative structure is governed jointly by coordinate geometry and generator curvature (commutator size).

\subsection{Cayley approximation error}
We characterize $\Cay(S)$ as a rational approximation to $\exp(S)$ and provide error bounds when $\norm{S}$ is small.
(We include a complete statement and proof in Appendix~\ref{app:cayley}.)


\section{Efficient Implementation}
\label{sec:efficiency}
\subsection{Low-rank Cayley application via Woodbury}
Let $S=UCU^\top$ with $C^\top=-C$ and $U\in\R^{d\times r}$.
To apply $x\mapsto \Cay(S)x$, we compute
\begin{equation}
  y \;=\; (I+\tfrac12 S)x \;=\; x + \tfrac12 U C (U^\top x),
\end{equation}
then solve $(I-\tfrac12 S)\,z = y$.
Using $z=y+U\alpha$ and multiplying by $(I-\tfrac12 S)$ gives a small $r\times r$ linear system for $\alpha$:
\begin{equation}
  \Big(I_r - \tfrac12 C (U^\top U)\Big)\alpha \;=\; \tfrac12 C (U^\top y),
\end{equation}
which we solve in $O(r^3)$, yielding overall $O(dr + r^3)$ cost per application.

\subsection{Differential rotor accumulation (DRA)}
For 1D positions with unit step $\Delta$, define a step generator $G(\Delta)$ and step rotor $R_\Delta=\Cay(G(\Delta))$.
Accumulate rotors as $R_{t+1}=R_\Delta R_t$ to avoid per-token generator evaluation and amortize computation.

\subsection{Kernel fusion}
We outline two integration modes:
(i) pre-rotate $Q,K$ before fused attention, and (ii) hybrid fusion where an Abelian core uses RoPE-like relative fusion and a low-rank coupling term is applied as a correction.
We provide JAX reference code and Triton pseudocode in the repository.


\section{Experiments}
\subsection{Benchmarks}
We evaluate on three settings:
\begin{enumerate}
  \item \textbf{Long-context text:} probe $\Delta$-extrapolation and ``needle-in-haystack'' retrieval performance.
  \item \textbf{2D/3D:} vision/video tasks where axis coupling is important.
  \item \textbf{Time--frequency audio:} spectrogram transformers where coupling time and frequency improves structure.
\end{enumerate}

\subsection{Ablations}
We sweep:
(i) rank $r\in\{0,4,8,16\}$ (with $r=0$ reducing to an Abelian baseline),
(ii) commutator budget strength (hard constraint or penalty),
(iii) Cayley vs exponential (when feasible),
(iv) differential rotor accumulation on/off,
(v) runtime overhead vs FlashAttention-style baselines.

\subsection{Diagnostics}
We log:
tokens/s, achieved TFLOPs, attention-kernel time, norm preservation error, and a \emph{drift diagnostic} measuring deviation from $\Delta$-only attention structure under increasing $\kappa$.


\section{Discussion}
BNR-PE provides a tunable path between strict relative positional structure and richer non-Abelian geometry.
The commutator budget behaves like a curvature control, suggesting a gauge-theoretic interpretation where non-commuting generators correspond to non-zero curvature.
Future work includes learned coordinate charts, path-ordered accumulation on grids, and multi-modal rotor alignment objectives.


\bibliographystyle{plain}
\bibliography{bnrpe}

\appendix
\section{Additional proofs and details}

\subsection{Cayley approximation bounds}
\label{app:cayley}
We recall that for skew $S$, $\Cay(S)$ is orthogonal and rationally approximates $\exp(S)$.
A standard expansion yields $\Cay(S)=\exp(S)+O(\norm{S}^3)$ as $\norm{S}\to 0$.
We provide explicit constants for operator-norm bounds under spectral radius assumptions.

\subsection{Implementation notes}
We include pseudocode for Woodbury-Cayley application and a note on numerical conditioning of the small $r\times r$ solve.


\end{document}
